% ============================================================
% Project: THE PHILOSOPHY OF VIABILITY
% Book: Book I — Ontology of Viability
% Book ID: VS.B01
% Part: Part I — Identity and Manifestation
% Chapter: Chapter 07 — Formal and Logical Reconstruction
% File: 4_ASSETS/diagrams/VS.B01.P01.C07.S02.SS02.P001.DG01.levels_of_ontological_and_formal_articulation.tex
% Asset: DG01 — Levels of Ontological and Formal Articulation
% Purpose: To distinguish ontological subject matter,
% conceptual articulation, and formal representation.
%
% Requires (preamble):
%   \usepackage{tikz}
%   \usetikzlibrary{arrows.meta, positioning}
% ============================================================

\begin{figure}[htbp]
    \centering
    \hyphenpenalty=10000 \exhyphenpenalty=10000\relax  % no hyphenation inside diagram labels

    % separator used inside the level boxes
    \def\vsdot{\enspace$\cdot$\enspace}

    \begin{tikzpicture}[
        font=\small,
        node distance=1.6cm,
        % --- base box ---
        box/.style={
            draw, rounded corners=3pt, align=center,
            text width=8.0cm, minimum height=1.9cm,
            inner sep=7pt, line width=0.8pt
        },
        % --- node roles ---
        ontbox/.style={box, fill=blue!8,  draw=blue!55!black},
        concbox/.style={box, fill=blue!16, draw=blue!55!black},
        formbox/.style={box, fill=gray!10, draw=black!60},
        % --- arrows (methodological relations, not production) ---
        directed/.style={-{Stealth[length=2.5mm]}, line width=0.9pt},
        feedback/.style={-{Stealth[length=2.5mm]}, line width=0.9pt, dashed},
        elabel/.style={font=\footnotesize\itshape, text=black!75,
                       align=left, text width=3.0cm}
    ]

    % ---------------------------------------------------------------
    % Three levels
    % ---------------------------------------------------------------
    \node[ontbox] (ontological) {
        \textbf{Ontological Subject Matter}\\[3pt]
        Being\vsdot Existent\vsdot Identity\vsdot Determination\\
        Relations, conditions, and manifestations
    };
    \node[concbox, below=of ontological] (conceptual) {
        \textbf{Conceptual Articulation}\\[3pt]
        Terms\vsdot Distinctions\vsdot Definitions\vsdot Claims\\
        Interpretations and ontological arguments
    };
    \node[formbox, below=of conceptual] (formal) {
        \textbf{Formal Representation}\\[3pt]
        Variables\vsdot Predicates\vsdot Relations\vsdot Operators\\
        Rules, models, and inferential structures
    };

    % ---------------------------------------------------------------
    % Articulation (solid, downward, left) and critical feedback
    % (dashed, upward, centre-right); labels sit beside the arrows
    % ---------------------------------------------------------------
    \draw[directed] ([xshift=-3.3cm]ontological.south) --
        node[elabel, right=5pt] {is articulated as a subject of inquiry}
        ([xshift=-3.3cm]conceptual.north);
    \draw[directed] ([xshift=-3.3cm]conceptual.south) --
        node[elabel, right=5pt] {is represented through formal expressions}
        ([xshift=-3.3cm]formal.north);

    \draw[feedback] ([xshift=0.8cm]conceptual.north) --
        node[elabel, right=5pt] {ontological adequacy constrains conceptual articulation}
        ([xshift=0.8cm]ontological.south);
    \draw[feedback] ([xshift=0.8cm]formal.north) --
        node[elabel, right=5pt] {formal analysis may require conceptual revision}
        ([xshift=0.8cm]conceptual.south);

    % ---------------------------------------------------------------
    % Legend
    % ---------------------------------------------------------------
    \node[below=0.6cm of formal, font=\scriptsize] {
        \tikz[baseline=-0.5ex]\draw[directed] (0,0) -- (0.7,0);~articulation
        \quad\quad
        \tikz[baseline=-0.5ex]\draw[feedback] (0,0) -- (0.7,0);~critical feedback
    };

    \end{tikzpicture}

    \caption{The distinction between ontological subject matter, conceptual articulation, and formal representation. The arrows indicate methodological relations and critical feedback, not the production of ontological reality by concepts or formal systems.}
    \label{fig:ch07_levels_of_articulation}
\end{figure}

% ============================================================
% END OF ASSET DG01 — CHAPTER 07, SECTION 02
% ============================================================
